\chapter*{Introducción}
\markboth{Introducción}{}
\addcontentsline{toc}{chapter}{Introducción}

     
La introducción es lo último que se escribe, pues en ella se va a dejar registro de lo elaborado en varios de los capítulos del desarrollo del proyecto. La introducción suele iniciar describiendo brevemente el problema que se resolvió. A continuación, se menciona también de forma breve cómo se solucionó el problema, también se presentan los resultados obtenidos.  

En principio se puede utilizar la introducción presentada en el anteproyecto, sin embargo, esta se deberá revisar al finalizar el proyecto. Tenga en cuenta que esta introducción serpa leída cuando el proyecto este finalizado por lo que se deberá escribir en pasado (o presente en los casos que aplique). Ej. Se desarrolló un …, El modelo realizado tuvo un desempeño .. , El sistema presenta


Se debe usar referencias bibliográficas en esta sección, por ejemplo:

La estimación formal de SoW y SioW se realiza mediante un modelo jerárquico Poisson--Binomial (\emph{pogit})\cite{glady2009predicting}
, el cual permite inferir simultáneamente la probabilidad de asignación del gasto a la empresa focal y la intensidad latente de consumo del cliente en la categoría, aun cuando el gasto total del cliente no sea observable. Dado el carácter latente de estas magnitudes, la evaluación del modelo no se plantea como un ejercicio de predicción puntual tradicional, sino como un análisis de coherencia estructural, estabilidad e interpretabilidad económica de las estimaciones obtenidas fuera de muestra.

